%% Chapter 9, Section 9.5 — Exercises
%% A First Course in Monte Carlo Methods
%% Sanz-Alonso & Al-Ghattas

\section{Exercises}
\label{sec:9.5}

\begin{enumerate}

\item[9.1]
\begin{enumerate}[label=\alph*)]
  \item Recall from Chapter~3 that the effective sample size
    \[
      \mathrm{ESS}_N = \frac{1}{\displaystyle\sum_{n=1}^{N} w_n^2},
    \]
    where $\{w_n\}_{n=1}^{N}$ are normalized weights, is often used to monitor the
    performance of importance sampling. Since particle filters are based on importance
    sampling, this quantity is also used to monitor particle filters. Show that
    $\mathrm{ESS}_N$ takes values between $1$ and $N$.

  \item A similar assessment is provided by the \emph{perplexity}, defined by
    \[
      \mathrm{perplex}_N := \frac{1}{N}\exp\!\left(-\sum_{n=1}^{N} w_n \log(w_n)\right),
    \]
    where again $\{w_n\}_{n=1}^{N}$ are normalized weights. Show that the perplexity
    provides an estimate of $\exp\!\left(-d_{\mathrm{KL}}(f \| g)\right)$, where
    \[
      d_{\mathrm{KL}}(f \| g) := \int \log\!\left(\frac{f(x)}{g(x)}\right) f(x)\,dx
    \]
    is the Kullback--Leibler divergence between $f$ and $g$. Show that
    $\mathrm{perplex}_N$ takes values between $1/N$ and $1$.
\end{enumerate}

\item[9.2] In this exercise, you will derive the Kalman filter, an important algorithm
that characterizes the filtering distributions for linear-Gaussian hidden Markov models.
Consider the following setting:
\begin{alignat*}{3}
  \text{(Initial condition)}  &\quad& X_0 &\sim \Normal(m_0, C_0), \\
  \text{(Dynamics model)}     &\quad& X_{j+1} &= MX_j + \xi_j, &\quad
    \xi_j &\sim \Normal(0, \Sigma), \\
  \text{(Observation model)}  &\quad& Y_{j+1} &= HX_{j+1} + \eta_{j+1}, &\quad
    \eta_j &\sim \Normal(0, \Gamma),
\end{alignat*}
where $X_0$, $\{\xi_j\}$, and $\{\eta_j\}$ are all assumed to be independent. Here,
$M \in \R^{d \times d}$ and $H \in \R^{k \times d}$ are given matrices that describe,
respectively, linear dynamics and linear observation models. Let $\widehat{f}_{j+1}$
denote the distribution of $X_{j+1} | Y_{1:j}$ and let $f_{j+1}$ be the filtering
distribution, i.e.\ the distribution of $X_{j+1} | Y_{1:j+1}$.

\begin{enumerate}[label=\alph*)]
  \item Show that $\widehat{f}_{j+1}$ and $f_{j+1}$ are Gaussian. Denote by
    $\widehat{m}_{j+1}$ and $\widehat{C}_{j+1}$ the mean and covariance of
    $\widehat{f}_{j+1}$. Likewise, denote by $m_{j+1}$ and $C_{j+1}$ the mean and
    covariance of $f_{j+1}$. In the next two parts of this question, you will find
    recursive formulas for these means and covariances.

  \item Show that
    \begin{align*}
      \widehat{m}_{j+1} &= Mm_j, \\
      \widehat{C}_{j+1} &= MC_j M^\top + \Sigma.
    \end{align*}

  \item Show that
    \begin{align*}
      C_{j+1}^{-1} &= \left(MC_j M^\top + \Sigma\right)^{-1} + H^\top \Gamma^{-1} H, \\
      C_{j+1}^{-1} m_{j+1}
        &= \left(MC_j M^\top + \Sigma\right)^{-1} Mm_j + H^\top \Gamma^{-1} Y_{j+1}.
    \end{align*}

  \item Using Woodbury's matrix identity, show that the mean $m_{j+1}$ and covariance
    $C_{j+1}$ of the filtering distribution can equivalently be written as
    \begin{align*}
      m_{j+1} &= \widehat{m}_{j+1} + K_{j+1} d_{j+1}, \\
      C_{j+1} &= (I - K_{j+1} H)\widehat{C}_{j+1},
    \end{align*}
    where
    \begin{align*}
      d_{j+1} &= Y_{j+1} - H\widehat{m}_{j+1}, \\
      S_{j+1} &= H\widehat{C}_{j+1} H^\top + \Gamma, \\
      K_{j+1} &= \widehat{C}_{j+1} H^\top S_{j+1}^{-1}.
    \end{align*}
    These formulas define the Kalman filter.
\end{enumerate}

\item[9.3] Consider the linear-Gaussian hidden Markov model
\begin{alignat*}{3}
  \text{(Initial condition)}  &\quad& X_0 &\sim \Normal(m_0, C_0), \\
  \text{(Dynamics model)}     &\quad& X_{j+1} &= MX_j + U_j + \xi_j, &\quad
    \xi_j &\sim \Normal(0, 0.01I), \\
  \text{(Observation model)}  &\quad& Y_{j+1} &= HX_{j+1} + \eta_j, &\quad
    \eta_j &\sim \Normal(0, 0.01),
\end{alignat*}
where $X_0$, $\{\xi_j\}$, and $\{\eta_j\}$ are all assumed to be independent. Set
\[
  M = \begin{bmatrix} \cos(6) & \sin(6) \\ -\sin(6) & \cos(6) \end{bmatrix},
  \quad
  H = \begin{bmatrix} 0.5 & 0.25 \end{bmatrix},
  \quad
  m_0 = \begin{bmatrix} 1 \\ 1 \end{bmatrix},
  \quad
  C_0 = \frac{1}{4}\begin{bmatrix} 7 & 5 \\ 5 & 7 \end{bmatrix},
\]
\[
  U_0 = \begin{bmatrix} 7 \\ 2 \end{bmatrix},\quad
  U_1 = \begin{bmatrix} 5 \\ 5 \end{bmatrix},\quad
  U_2 = \begin{bmatrix} -1 \\ 2 \end{bmatrix},\quad
  U_3 = \begin{bmatrix} -1 \\ -2 \end{bmatrix},\quad
  U_4 = \begin{bmatrix} 1 \\ -3 \end{bmatrix}.
\]
Note that $X_j \in \R^2$ and $Y_j \in \R$.

\begin{enumerate}[label=\alph*)]
  \item Generate artificial data $Y_1, \ldots, Y_5$. To that end, draw $X_0$ from the
    initial condition, draw $X_1, \ldots, X_5$ from the dynamics model, and draw data
    $Y_1, \ldots, Y_5$ from the observation model. Use these data for the next two parts
    of this question.

  \item Implement a bootstrap particle filter with $N = 1000$ particles. Plot the
    particle locations versus time for time $j = 0, \ldots, 5$. For each time, find
    explicitly the conditional distribution of $X_j | Y_1, \ldots, Y_j$ and plot (on top
    of the previous plots) the $95\%$ contour of the conditional density.

  \item Implement the optimal particle filter with the same specifications as before.
    Compare the results.
\end{enumerate}

\item[9.4] Consider the Lorenz~63 model
\begin{align*}
  \frac{dv}{dt} &= \sigma(w - v), \\
  \frac{dw}{dt} &= v(\rho - z) - w, \\
  \frac{dz}{dt} &= vw - \beta z,
\end{align*}
with the standard parameter values $\sigma = 10$, $\rho = 28$, $\beta = \frac{8}{3}$.
For given $\Delta t > 0$, let $a : \R^3 \to \R^3$ denote the $\Delta t$-flow, so that,
for $x = (v, w, z) \in \R^3$, $a(x)$ denotes the solution of the Lorenz~63 equation at
time $\Delta t$ with initial condition $x$. In this exercise, you will implement a
particle filter for a stochastic Lorenz~63 model which is observed at discrete times
$j\Delta t$, $1 \leq j \leq J$, for a time window of length $J = 100$. Here, $\Delta t$
represents the time-span between observations, which will be assumed to be constant.
Precisely, we consider the following hidden Markov model:
\begin{alignat*}{3}
  \text{(Initial condition)}  &\quad& X_0 &\sim \Normal(m_0, C_0), \\
  \text{(Dynamics model)}     &\quad& X_{j+1} &= a(X_j) + \xi_j, &\quad
    \xi_j &\sim \Normal(0, \Sigma), \\
  \text{(Observation model)}  &\quad& Y_{j+1} &= X_{j+1} + \eta_j, &\quad
    \eta_j &\sim \Normal(0, \Gamma),
\end{alignat*}
where $X_0$, $\{\xi_j\}$, and $\{\eta_j\}$ are all assumed to be independent. Set
$m_0 = (0, 0, 0)$ and $C_0 = \Sigma = \Gamma = 10^{-4} I$.

\begin{enumerate}[label=\alph*)]
  \item Draw an initial condition $X_0 \sim \Normal(m_0, C_0)$ and simulate the dynamics
    model to obtain a true state trajectory $\{X_j\}_{j=1}^{J}$. Next, using the
    observation model, simulate data $\{Y_j\}_{j=1}^{J}$. Your goal will be to recover
    the true signal $\{X_j\}_{j=1}^{J}$ from the simulated data $\{Y_j\}_{j=1}^{J}$.

  \item Implement a bootstrap particle filter with $N = 100$ particles. Plot together the
    true state trajectory, the particle filter estimates, and the observations for each
    time step. Your plot should comprise three subplots, corresponding to each component
    of the Lorenz system.

  \item Implement the optimal particle filter with the same specifications as before.
    Compare the results.
\end{enumerate}

\end{enumerate}
